\documentclass[11pt]{article}
\usepackage{amsmath,amssymb,amsthm}
\usepackage[margin=1in]{geometry}

\newtheorem{theorem}{Theorem}
\newtheorem{lemma}[theorem]{Lemma}
\newtheorem{definition}[theorem]{Definition}
\newtheorem{proposition}[theorem]{Proposition}
\newtheorem{remark}[theorem]{Remark}
\newtheorem{corollary}[theorem]{Corollary}

\title{Semantic Soundness of the Well-Chained Fragment:\\
Nelson's Program in Guard--Rose Recursive Arithmetic}
\author{Summary of Agda development}
\date{April 2026}

\begin{document}
\maketitle

\section{Introduction}

Building on the internal verification experiments of the previous
session, we investigate the semantic soundness of the proof checker
$\mathit{thFunTFor}$. The main results are:

\begin{enumerate}
\item A modular dispatch/extractor library for Nelson-style branch
  verification (Assembly~1 infrastructure).
\item A corrected semantic predicate: $\mathit{TreeEq}(L,R)=O$ is the
  \emph{wrong} predicate; a ground-term evaluator $\mathit{evalCode}$
  gives the right one.
\item An inductive well-chained fragment with semantic soundness
  proved by induction on the fragment structure.
\item A full Nelson pipeline: from a specific $\mathit{Deriv}$ proof,
  through encoding, through well-chainedness, to internal semantic
  soundness.
\end{enumerate}

All results are formalized in Agda with \texttt{--safe --without-K
--exact-split}. No postulates are used. All files compile in under
0.2 seconds.

\section{Modular Branch Verification Library}

\subsection{NelsonDispatch}

We factor the 26-branch dispatch chain into a reusable library.
For each tag $k \in \{0,\ldots,25\}$, $\mathit{ndDispatchToCase}_k$
proves:
\[
  \mathit{ndDispatch}(\mathit{Pair}(\mathit{reify}(\mathit{natCode}\;k),
  \mathit{data}'), \mathit{recs})
  = \mathit{case}_k(\ldots)
\]
using $\mathit{ndBranchMiss}$ and $\mathit{ndBranchHit}$ from
$\mathit{ThFunTForCorrectDefs}$. The inequality proofs are all
$\mathit{refl}$ (since $\mathit{natEq}\;k\;n = \mathit{false}$ for
distinct concrete naturals).

Also provides $\mathit{phase1Nd}$: the combined
$\mathit{axRecNd} + \mathit{guardNd}$ step.

\subsection{NelsonExtractors}

Shared reduction lemmas:
\begin{itemize}
\item $\mathit{origARed}$, $\mathit{origBRed}$: extract components
  from $\mathit{orig} = \mathit{Pair}(\mathit{tag}, \mathit{Pair}(x,y))$.
\item $\mathit{innerPassthrough}$: the full chain
  $\mathit{axRecNd} + \mathit{guardNd} + \mathit{ndDispatchPairMiss}
  + \mathit{sndArg2}$ for Pair-Pair-tagged sub-encodings.
\item $\mathit{recsBWithPassthrough}$, $\mathit{recsBLWithPassthrough}$,
  $\mathit{recsBRWithPassthrough}$: recursive result extraction
  via passthrough.
\item $\mathit{mkAp1FRed}$, $\mathit{mkAp2FRed}$, $\mathit{mkEqFRed}$:
  combinator reduction helpers.
\end{itemize}

\subsection{New Case Proofs}

Using the library, three new cases are proved:
\begin{center}
\begin{tabular}{llll}
\hline
\textbf{File} & \textbf{Tag} & \textbf{Rule} & \textbf{Type} \\
\hline
NelsonAxI & 0 & axI & Base \\
NelsonAxFst & 1 & axFst & Base \\
NelsonRuleSym & 18 & ruleSym & Unary \\
\hline
\end{tabular}
\end{center}

Each file is 55--65 lines. The library makes adding cases mechanical.

\section{Experiment E: The Wrong and Right Predicates}

\subsection{The Wrong Predicate}

Define:
\[
  \mathit{Good}(p) := \mathit{TreeEq}(\mathit{Fst}(\mathit{thFunTFor}(p)),\;
  \mathit{Snd}(\mathit{thFunTFor}(p))) = O
\]

\begin{theorem}[NelsonExpE: Failure]
For $p = \mathit{enc}_{\mathit{axI}}(\mathit{code}(O))$:
\[
  \mathit{TreeEq}(\mathit{Fst}(\mathit{thFunTFor}(p)),\;
  \mathit{Snd}(\mathit{thFunTFor}(p))) = \mathit{Pair}(O,O)
\]
is derivable in $\mathit{Deriv}$. That is, $\mathit{Good}(p)$
gives $\mathit{Pair}(O,O)$, not~$O$.
\end{theorem}

The reason: $\mathit{thFunTFor}$ returns $\mathit{Pair}(L,R)$
where $L = \mathit{code}(I(O))$ and $R = \mathit{code}(O)$. These
are \emph{different tree structures}: $L = \mathit{mkAp1}(\mathit{codeI},
\mathit{code}(O))$ while $R = \mathit{code}(O)$. The equation
$I(O) = O$ is \emph{true}, but the \emph{codes} differ structurally.
$\mathit{TreeEq}$ compares trees, not semantic content.

\subsection{The Right Predicate}

Define $\mathit{evalCode} : \mathit{Fun1} = \mathit{Rec}\;O\;
\mathit{evalStep}$, a ground-term evaluator on coded terms.
The step function dispatches on the term tag:
\begin{itemize}
\item $\mathit{tagO}$: return~$O$.
\item $\mathit{tagAp1}$: return $\mathit{evalCode}$ of the argument
  code (strips the functor application; correct for~$I$).
\item Default: passthrough ($\mathit{sndArg2}$).
\end{itemize}

Then define:
\[
  \mathit{EqTrue}(p) := \mathit{TreeEq}(\mathit{evalCode}(
  \mathit{Fst}(\mathit{thFunTFor}(p))),\;
  \mathit{evalCode}(\mathit{Snd}(\mathit{thFunTFor}(p)))) = O
\]

\begin{theorem}[EvalCodeCorrect: Success]
\label{thm:eqtrue}
For $p = \mathit{enc}_{\mathit{axI}}(\mathit{code}(O))$:
\[
  \mathit{EqTrue}(p)
\]
is derivable in $\mathit{Deriv}$. Concretely:
$\mathit{evalCode}(\mathit{code}(I(O))) = O$ and
$\mathit{evalCode}(\mathit{code}(O)) = O$,
so $\mathit{TreeEq}(O, O) = O$.
\end{theorem}

\section{Semantic Soundness per Rule}

\begin{definition}[SemTrue]
$\mathit{SemTrue}(p) := \mathit{Deriv}\;\mathit{hyp}\;
(\mathit{eqn}\;(\mathit{evalCode}(\mathit{Fst}(\mathit{thFunTFor}(p))))
\;(\mathit{evalCode}(\mathit{Snd}(\mathit{thFunTFor}(p)))))$.
\end{definition}

\begin{theorem}[SemSoundRefl]
$\mathit{SemTrue}(\mathit{enc}_{\mathit{axRefl}}(x))$ holds
unconditionally.
\end{theorem}

\begin{theorem}[SemSoundAxI]
$\mathit{SemTrue}(\mathit{enc}_{\mathit{axI}}(\mathit{code}(O)))$
holds unconditionally (ground instance).
\end{theorem}

\begin{theorem}[SemSoundSym]
If $\mathit{SemTrue}(\mathit{sp})$, then
$\mathit{SemTrue}(\mathit{enc}_{\mathit{ruleSym}}(\mathit{sp}))$.
No side condition.
\end{theorem}

\begin{theorem}[SemSoundTrans]
If $\mathit{SemTrue}(\mathit{sp}_1)$ and
$\mathit{SemTrue}(\mathit{sp}_2)$ and
\[
  \mathit{Snd}(\mathit{thFunTFor}(\mathit{sp}_1)) =
  \mathit{Fst}(\mathit{thFunTFor}(\mathit{sp}_2))
  \quad\text{(middle term match)}
\]
then $\mathit{SemTrue}(\mathit{enc}_{\mathit{ruleTrans}}(\mathit{sp}_1,
\mathit{sp}_2))$.
\end{theorem}

The middle term match is the \emph{only} side condition.
It holds automatically for valid proof encodings (both sides
are $\mathit{code}(u)$ for the shared middle term~$u$), but
fails for arbitrary trees.

\section{The Well-Chained Fragment}

\begin{definition}[WC]
The inductive type $\mathit{WC}\;p\;\mathit{hyp}$ characterizes
well-chained proof encodings:
\begin{itemize}
\item $\mathit{wcRefl}$: $\mathit{WC}(\mathit{enc}_{\mathit{axRefl}}(x))$.
\item $\mathit{wcSym}$: $\mathit{WC}(\mathit{sp}) \to
  \mathit{WC}(\mathit{enc}_{\mathit{ruleSym}}(\mathit{sp}))$.
\item $\mathit{wcTrans}$: $\mathit{WC}(\mathit{sp}_1) \to
  \mathit{WC}(\mathit{sp}_2) \to \mathit{match} \to
  \mathit{WC}(\mathit{enc}_{\mathit{ruleTrans}}(\mathit{sp}_1,
  \mathit{sp}_2))$.
\end{itemize}
\end{definition}

\begin{theorem}[Semantic Soundness of the Well-Chained Fragment]
\label{thm:semsound}
For any $p$ with $\mathit{WC}\;p\;\mathit{hyp}$:
\[
  \mathit{SemTrue}(p)
\]
Proof by induction on the $\mathit{WC}$ derivation.
\end{theorem}

\section{The Nelson Pipeline}

For a specific proof $\mathit{Trans}(\mathit{Refl}(x),
\mathit{Refl}(x))$, the full pipeline is:

\begin{enumerate}
\item \textbf{Encode}: $\mathit{encRefl} = \mathit{Pair}(\mathit{tag}_{17},
  \mathit{Pair}(x, O))$.
\item \textbf{Decode}: $\mathit{nelsonAxRefl}$ gives
  $\mathit{thFunTFor}(\mathit{encRefl}) = \mathit{Pair}(x, x)$.
\item \textbf{Match}: $\mathit{Snd}(\mathit{Pair}(x,x)) = x =
  \mathit{Fst}(\mathit{Pair}(x,x))$ --- trivially.
\item \textbf{WC}: $\mathit{wcTrans}(\mathit{wcRefl}, \mathit{wcRefl},
  \mathit{match})$.
\item \textbf{SemSound}: $\mathit{semSound}$ gives
  $\mathit{evalCode}(L) = \mathit{evalCode}(R)$.
\item \textbf{TreeEq}: $\mathit{TreeEq}(\mathit{evalCode}(L),
  \mathit{evalCode}(R)) = O$.
\end{enumerate}

A deeper example, $\mathit{Sym}(\mathit{Trans}(\mathit{Refl}(x),
\mathit{Refl}(x)))$, is also proved, demonstrating composition
through three rules.

\begin{remark}[Connection to Nelson]
This pipeline formalizes Nelson's ``for each specific proof, the
bounds hold'' strategy from \emph{Elements}~\cite{nelson-elements}.
The matching conditions are \emph{derivable consequences} of
$\mathit{thFunTFor}$ correctness for each specific proof.
No universal claim over all trees is needed.
\end{remark}

\section{The Boundary}

\subsection{Three Obstructions to Universal $\mathit{Good}(p)$}

\begin{enumerate}
\item \textbf{Base case}: $\mathit{thFunTFor}(\mathit{lf}) = O$,
  and $\mathit{Fst}(O)$/$\mathit{Snd}(O)$ are stuck.
\item \textbf{Middle term match}: $\mathit{ruleTrans}$ needs
  $R_1 = L_2$, which fails for arbitrary trees (invalid proof trees exist).
\item \textbf{TreeEq inversion}: even with the match, going from
  $\mathit{TreeEq}(\mathit{evalCode}(L), \mathit{evalCode}(R)) = O$
  to the equational fact $\mathit{evalCode}(L) = \mathit{evalCode}(R)$
  requires inverting $\mathit{TreeEq}$.
\end{enumerate}

\subsection{The Precise Boundary}

The obstruction is \textbf{not}:
\begin{itemize}
\item semantic composition (it works),
\item recursion or induction (WC induction works),
\item proof-checker complexity or size blow-up (all $< 0.2$s).
\end{itemize}

The obstruction \textbf{is}: the validity side conditions (e.g.,
the middle term match for $\mathit{ruleTrans}$) are not uniformly
provable for arbitrary proof trees. Validity is a \emph{semantic
invariant extracted by the checker itself}, creating a dependency
loop between validity and computation.

This is the precise point where Nelson's program meets G\"odel's
second theorem: for each specific proof, the conditions hold; but
the universal claim requires the finitary credo, which Nelson
considers circular.

\section{File Inventory}

\begin{center}
\begin{tabular}{lrl}
\hline
\textbf{File} & \textbf{Lines} & \textbf{Role} \\
\hline
NelsonDispatch.agda & 317 & Dispatch chains for all 26 tags \\
NelsonExtractors.agda & 339 & Shared extractors + passthrough \\
NelsonAxI.agda & 64 & Tag 0 base case \\
NelsonAxFst.agda & 62 & Tag 1 base case \\
NelsonRuleSym.agda & 66 & Tag 18 unary case \\
NelsonExpE.agda & 67 & Experiment E failure \\
EvalCode.agda & 78 & Ground-term evaluator \\
EvalCodeCorrect.agda & 231 & evalCode correctness \\
EvalCodeTrans.agda & 65 & Ground composition lemma \\
SemSoundRefl.agda & 54 & Semantic soundness: axRefl \\
SemSoundAxI.agda & 69 & Semantic soundness: axI \\
SemSoundSym.agda & 96 & Semantic soundness: ruleSym \\
SemSoundTrans.agda & 119 & Semantic soundness: ruleTrans \\
WellChained.agda & 130 & WC fragment + inductive soundness \\
NelsonPipeline.agda & 157 & Full pipeline demonstration \\
\hline
\textbf{Total} & \textbf{1914} & \\
\hline
\end{tabular}
\end{center}

All files compile in under 0.2 seconds. No postulates.

\begin{thebibliography}{9}
\bibitem{rose1961} H.\,E.\ Rose,
\emph{On the Consistency and Undecidability of Recursive Arithmetic},
Zeitschr.\ f.\ math.\ Logik und Grundlagen d.\ Math., Bd.~7,
S.~124--135 (1961).

\bibitem{guard1963} J.\,R.\ Guard,
\emph{Lecture Notes on Recursive Arithmetic},
Princeton University, Mathematics Department (1963).

\bibitem{nelson} E.\ Nelson,
\emph{Predicative Arithmetic},
Mathematical Notes~32, Princeton University Press (1986).

\bibitem{nelson-elements} E.\ Nelson,
\emph{Elements},
arXiv:1510.00369 [math.LO] (2015), work in progress.
\end{thebibliography}

\end{document}
